\documentclass[aspectratio=169,11pt]{beamer}

\usetheme{Madrid}
\usecolortheme{default}
\usepackage[utf8]{inputenc}
\usepackage[T1]{fontenc}
\usepackage{amsmath,amssymb}
\usepackage{graphicx}
\usepackage{booktabs}
\usepackage{array}
\usepackage{xcolor}

\definecolor{frcblue}{RGB}{28,45,68}
\definecolor{frcaccent}{RGB}{18,118,150}
\setbeamercolor{structure}{fg=frcblue}
\setbeamercolor{frametitle}{fg=white,bg=frcblue}
\setbeamertemplate{headline}{}
\setbeamertemplate{navigation symbols}{}
\setbeamertemplate{footline}[frame number]
\setbeamertemplate{frametitle}{%
  \nointerlineskip%
  \begin{beamercolorbox}[wd=\paperwidth,ht=0.62cm,dp=0.22cm,leftskip=0.45cm,rightskip=0.45cm]{frametitle}
    \usebeamerfont{frametitle}\insertframetitle
  \end{beamercolorbox}%
}

\title[Advanced Flying Robots -- INDI Controller]{Advanced Flying Robots -- INDI Controller}
\subtitle{Current Progress CW4}
\author{Georg Wolnik}
\date{Milestone 1: Geometrical Controller \\ Milestone 2: Motion Planning}

\newcommand{\img}[2]{%
  \begin{center}
    \includegraphics[width=#1\linewidth,height=0.72\textheight,keepaspectratio]{#2}
  \end{center}
}

\newcommand{\twopics}[4]{%
  \begin{columns}[c,totalwidth=\textwidth]
    \begin{column}{0.5\textwidth}
      \centering
      \includegraphics[width=\linewidth,height=0.76\textheight,keepaspectratio]{#2}
    \end{column}
    \begin{column}{0.5\textwidth}
      \centering
      \includegraphics[width=\linewidth,height=0.76\textheight,keepaspectratio]{#4}
    \end{column}
  \end{columns}
}

\begin{document}

\begin{frame}
  \titlepage
\end{frame}

\begin{frame}{Current Progress: Two Major Milestones}
\begin{block}{Milestone 1 --- Geometrical controller working on hardware}
\begin{itemize}
  \item Geometrical controller is running reliably enough for Figure-8 flight tests.
  \item Gain tuning improved tracking quality, reaching approximately \(8\,\mathrm{cm}\) RMSE in the best current Figure-8 result.
  \item Flight logs and analysis plots provide the main evidence for this milestone.
\end{itemize}
\end{block}

\begin{block}{Milestone 2 --- Motion planning pipeline}
\begin{itemize}
  \item Four planning modes are implemented and comparable: fixed-time spline, Richter timing, explicit SE(3), and paper-style time redistribution.
  \item The \(k_t\) scan functionality searches aggressive but feasible timing parameters.
  \item Reference, limits, state, kinematics, and 3D orientation plots visualize and validate the planned trajectories.
\end{itemize}
\end{block}
\end{frame}

\section{Hardware Progress}

\begin{frame}{Geometric Controller On Hardware: Initial State}
\framesubtitle{\texttt{figure8\_20260423\_145011.csv}}
\twopics{Main analysis}{Controls/logs/figure8_20260423_145011_analysis.png}
        {3D orientation}{Controls/logs/figure8_20260423_145011_3d_orientation.png}
\end{frame}

\begin{frame}{Milestone 1 Evidence: Full Analysis And 3D Trajectory}
\framesubtitle{\texttt{figure8\_onboard\_m0\_s1-0x\_r1\_20260506\_162613.csv}}
\twopics{Full analysis}{Controls/logs/figure8_onboard_m0_s1-0x_r1_20260506_162613_analysis.png}
        {3D trajectory}{Controls/logs/figure8_onboard_m0_s1-0x_r1_20260506_162613_3d_orientation.png}
\end{frame}

\begin{frame}{Milestone 1 Evidence: Kinematics And State Axes}
\twopics{Kinematics}{Controls/logs/figure8_onboard_m0_s1-0x_r1_20260506_162613_analysis_kinematics.png}
        {State axes}{Controls/logs/figure8_onboard_m0_s1-0x_r1_20260506_162613_analysis_axes.png}
\end{frame}

\section{Motion Planning}

\begin{frame}{Milestone 2 Overview: Motion Planning Modes}
\small
\begin{tabular}{p{0.08\textwidth}p{0.32\textwidth}p{0.20\textwidth}p{0.22\textwidth}}
\toprule
Mode & Core idea & Attitude source & Time scaling \\
\midrule
0 & Base spline trajectory & Flatness-derived & Manual / fixed \\
1 & Richter-like minimum-snap with time allocation & Flatness-derived & \(k_t\)-based \\
2 & Explicit SE(3): authored attitude + position & Explicit quaternion/SO(3) path & Richter timing + SO(3)/SLERP attitude \\
3 & Paper-style flat-output planning & Flatness-derived & Iterative time redistribution \\
\bottomrule
\end{tabular}

\vspace{0.4cm}
\begin{block}{Key point}
The modes are not four unrelated planners. They share a smooth-polynomial trajectory backbone; the important differences are segment timing and attitude generation.
\end{block}
\end{frame}

\begin{frame}{Planning Formulation: Polynomial And Snap Cost}
\begin{block}{Normalized segment polynomial}
\[
\hat t=\frac{\tau}{T_i},\qquad \hat t\in[0,1],
\qquad
\sigma_i(\hat t)=\sum_{k=0}^{8}c_{i,k}\hat t^k
\]
\end{block}

\begin{block}{Physical-time derivative scaling}
\[
\frac{d^r\sigma_i}{d\tau^r}
=
\frac{1}{T_i^r}
\frac{d^r\sigma_i}{d\hat t^r}
\]
\end{block}

\begin{block}{Per-axis snap cost}
\[
J_i
=
\frac{1}{T_i^7}
\int_0^1
\left(
\frac{d^4\sigma_i}{d\hat t^4}
\right)^2
d\hat t,
\qquad
J=\sum_i J_i
\]
\end{block}
\end{frame}

\begin{frame}{Planning Formulation: QP Form And Equality Constraints}
\begin{block}{Current QP form}
\[
\min_{\mathbf c}\frac{1}{2}\mathbf c^\top Q\mathbf c+\mathbf q^\top\mathbf c
\quad\text{s.t.}\quad
A_{eq}\mathbf c=b_{eq}
\]
\small Solved independently for \(x,y,z\) and yaw where applicable; coupling re-enters through flatness and feasibility checks.
\end{block}

\begin{block}{Waypoint interpolation}
\[
\sigma_i(0)=w_i,\qquad \sigma_i(1)=w_{i+1}
\]
\end{block}

\begin{block}{Derivative continuity at segment junctions}
\[
\frac{1}{T_i^r}\sigma_i^{(r)}(1)
=
\frac{1}{T_{i+1}^r}\sigma_{i+1}^{(r)}(0),
\qquad r=1,\dots,4
\]
\end{block}
\end{frame}

\begin{frame}{Planning Formulation: Boundary Conditions}
\begin{block}{Non-periodic rest-to-rest endpoint constraints}
\[
\frac{d^r\sigma}{d\tau^r}(0)=0,
\qquad
\frac{d^r\sigma}{d\tau^r}(T)=0,
\qquad r=1,\dots,4
\]
\end{block}

\begin{block}{Periodic-loop condition}
\[
\frac{d^r\sigma_{\mathrm{last}}}{d\tau^r}(T_{\mathrm{last}})
=
\frac{d^r\sigma_{\mathrm{first}}}{d\tau^r}(0),
\qquad r=1,\dots,4
\]
\end{block}

\begin{block}{Important limitation}
The current spline/Richter QP uses equality constraints only. Thrust, body-rate, torque, and motor limits are checked after planning, not enforced as QP inequalities.
\end{block}
\end{frame}

\begin{frame}{Planning Formulation: Differential Flatness}
\begin{block}{Force direction}
\[
f=m(a+g e_3),\qquad z_B=\frac{f}{\|f\|}
\]
\end{block}

\begin{block}{Attitude from desired yaw}
\[
x_C=
\begin{bmatrix}
\cos\psi\\
\sin\psi\\
0
\end{bmatrix},
\qquad
y_B=\frac{z_B\times x_C}{\|z_B\times x_C\|},
\qquad
x_B=y_B\times z_B
\]
\[
R=\begin{bmatrix}x_B & y_B & z_B\end{bmatrix}
\]
\end{block}

\small Modes 0/1/3 use this flatness path for attitude, thrust, angular-rate, and feasibility references.
\end{frame}

\begin{frame}{Planning Validation: Feasibility Is Post-Solve}
\begin{block}{Core utilization ratios}
\[
u_T=\frac{|T|}{T_{\max}},
\qquad
u_\omega=\frac{\|\omega\|}{\omega_{\max}},
\qquad
u_\tau=\frac{\|\tau\|_\infty}{\tau_{\max}}
\]
\[
u_{\max}=\max(u_T,u_\omega,u_\tau)
\]
\end{block}

\begin{itemize}
  \item The limits dashboard visualizes sampled feasibility after the polynomial trajectory has been solved.
  \item Additional checks include motor-mixing feasibility, per-motor usage, rate proxies, sustained violation duration, and battery-sag proxy.
  \item These are model/proxy checks for planning risk, not direct proof of hardware saturation behavior.
\end{itemize}
\end{frame}

\begin{frame}{Planning Modes: What Changes Between Modes?}
\small
\begin{tabular}{p{0.10\textwidth}p{0.25\textwidth}p{0.47\textwidth}}
\toprule
Mode & What stands out & Formulation / explanation \\
\midrule
0 & Fixed timing baseline & Segment times \(T_s\) are provided directly; the optimizer solves only for polynomial coefficients. \\
1 & \(k_t\)-based timing & \(v_{avg}=0.5+1.5\sqrt{k_t}\), \quad \(T_i=\max(\|p_{i+1}-p_i\|/v_{avg},0.15)\). \\
2 & Explicit attitude & Position uses Richter-like timing, but attitude waypoints use \(q(s)=\operatorname{SLERP}(q_i,q_{i+1},s)\). \\
3 & Time redistribution & Starts from Mode 1 timing and iterates \(T_i^{target}=\max(T_{total}J_i/\sum_jJ_j,T_{\min})\), then \(T_i\leftarrow T_i+0.5(T_i^{target}-T_i)\), with safety clamps on segment time. \\
\bottomrule
\end{tabular}
\end{frame}

\begin{frame}{Timing Selection: Latest Figure-8 \(k_t\) Scan}
\small
The scan selected the largest sampled \(k_t\) values that stayed below the configured planning-side utilization limit for the Figure-8 planners shown in the main evidence:
\[
\texttt{target\_utilization\_max}=0.99
\]

\vspace{0.2cm}
\begin{tabular}{cccc}
\toprule
Mode & Figure-8 \(k_t\) & Status & Interpretation \\
\midrule
1 & \(0.3107\) & below\_band & Highest safe sampled Richter timing value \\
3 & \(0.7116\) & below\_band & Highest safe sampled paper-style timing value \\
\bottomrule
\end{tabular}

\vspace{0.4cm}
\footnotesize
\textbf{Note:} Mode 2 is shown in the appendix with a vertical loop, because explicit attitude planning is most meaningful for inversion-capable manoeuvres.
\end{frame}

\section{Main Planning Evidence}

\begin{frame}{Main Planning Evidence: Mode 1 Overview}
\framesubtitle{Minimum-snap with automatic \(k_t\) timing}
\img{0.86}{flying_drone_stack/results/planning_sim/reference/mode1/figure8/overview.png}
\end{frame}

\begin{frame}{Main Planning Evidence: Mode 1 Limits And Orientation}
\twopics{Limits dashboard}{flying_drone_stack/results/planning_sim/reference/mode1/figure8/limits_dashboard.png}
        {3D trajectory + orientation}{flying_drone_stack/results/planning_sim/reference/mode1/figure8/trajectory_3d_orientation.png}
\end{frame}

\begin{frame}{Main Planning Evidence: Mode 1 Kinematics And State Axes}
\twopics{Kinematics}{flying_drone_stack/results/planning_sim/reference/mode1/figure8/kinematics_axes.png}
        {State axes}{flying_drone_stack/results/planning_sim/reference/mode1/figure8/state_axes.png}
\end{frame}

\begin{frame}{Main Planning Evidence: Mode 3 Overview}
\framesubtitle{Richter-family planner plus iterative time redistribution}
\img{0.86}{flying_drone_stack/results/planning_sim/reference/mode3/figure8/overview.png}
\end{frame}

\begin{frame}{Main Planning Evidence: Mode 3 Limits And Orientation}
\twopics{Limits dashboard}{flying_drone_stack/results/planning_sim/reference/mode3/figure8/limits_dashboard.png}
        {3D trajectory + orientation}{flying_drone_stack/results/planning_sim/reference/mode3/figure8/trajectory_3d_orientation.png}
\end{frame}

\begin{frame}{Main Planning Evidence: Mode 3 Kinematics And State Axes}
\twopics{Kinematics}{flying_drone_stack/results/planning_sim/reference/mode3/figure8/kinematics_axes.png}
        {State axes}{flying_drone_stack/results/planning_sim/reference/mode3/figure8/state_axes.png}
\end{frame}

\begin{frame}{Roadmap: Milestones 3 And 4}
\begin{block}{Milestone 3 --- Motion capture integration}
\begin{itemize}
  \item Switch to the mocap system and make it work with our full setup.
  \item Run the geometrical controller on the drone using mocap-based state feedback.
  \item Expected result: significantly improved tracking error compared with optical-flow-only state estimation.
\end{itemize}
\end{block}

\begin{block}{Milestone 4 --- INDI controller implementation and validation}
\begin{enumerate}
  \item Implement / integrate the INDI controller path.
  \item Validate behavior in simulation first.
  \item Tune gains, filters, and sensor/noise handling.
  \item Validate safely on real hardware, starting with simple maneuvers.
  \item Extend to trajectory tracking experiments and compare against the geometrical controller.
\end{enumerate}
\end{block}
\end{frame}

\appendix
\section{Appendix}

\begin{frame}{Appendix: Mode 0 Fixed-Time Baseline}
\framesubtitle{Reference baseline with manually selected segment times}
\img{0.86}{flying_drone_stack/results/planning_sim/reference/mode0/figure8/overview.png}
\end{frame}

\begin{frame}{Appendix: Mode 0 Limits And Orientation}
\twopics{Limits dashboard}{flying_drone_stack/results/planning_sim/reference/mode0/figure8/limits_dashboard.png}
        {3D trajectory + orientation}{flying_drone_stack/results/planning_sim/reference/mode0/figure8/trajectory_3d_orientation.png}
\end{frame}

\begin{frame}{Appendix: Mode 0 Kinematics And State Axes}
\twopics{Kinematics}{flying_drone_stack/results/planning_sim/reference/mode0/figure8/kinematics_axes.png}
        {State axes}{flying_drone_stack/results/planning_sim/reference/mode0/figure8/state_axes.png}
\end{frame}

\begin{frame}{Appendix: Mode 2 Vertical Loop Overview}
\framesubtitle{Explicit SE(3) attitude path for inversion-capable maneuvers}
\img{0.86}{flying_drone_stack/results/planning_sim/reference/mode2/loop/overview.png}
\end{frame}

\begin{frame}{Appendix: Mode 2 Vertical Loop Limits And Orientation}
\twopics{Limits dashboard}{flying_drone_stack/results/planning_sim/reference/mode2/loop/limits_dashboard.png}
        {3D trajectory + orientation}{flying_drone_stack/results/planning_sim/reference/mode2/loop/trajectory_3d_orientation.png}
\end{frame}

\begin{frame}{Appendix: Mode 2 Vertical Loop Kinematics And State Axes}
\twopics{Kinematics}{flying_drone_stack/results/planning_sim/reference/mode2/loop/kinematics_axes.png}
        {State axes}{flying_drone_stack/results/planning_sim/reference/mode2/loop/state_axes.png}
\end{frame}

\end{document}
